\documentclass[letterpaper,journal]{IEEEtran}
\usepackage{amsmath,amsfonts}
\usepackage{algorithmic}
\usepackage{algorithm}
\usepackage{array}
\usepackage[caption=false,font=normalsize,labelfont=sf,textfont=sf]{subfig}
\usepackage{textcomp}
\usepackage{stfloats}
\usepackage{url}
\usepackage{verbatim}
\usepackage{graphicx}
\usepackage{booktabs}
\usepackage{multirow}
\usepackage[T1]{fontenc}
\usepackage[utf8]{inputenc}
\usepackage[english]{babel}
\usepackage{csquotes}
\usepackage{cite}
\hyphenation{op-tical net-works semi-conduc-tor IEEE-Xplore}

\begin{document}

\title{A Computer-Vision-Based Bikeability Score}

\author{Bünyamin~Aydemir,
        Gonzalo~Guerrero~Salazar,
        and~Hendrick~Fischer% <-this % stops a space
\thanks{Manuscript created July 2026. The authors are students at the
Baden-Württemberg Cooperative State University (DHBW) within the
Computer Vision course.}}

% The paper headers
\markboth{Computer Vision -- Project Work, DHBW, 2026}%
{Aydemir \MakeLowercase{\textit{et al.}}: A Computer-Vision-Based Bikeability Score}

\maketitle

\begin{abstract}
% --- NOTE: This section will be written later (part of page 1). ---
\end{abstract}

\begin{IEEEkeywords}
Bikeability, computer vision, zero-shot classification, vision-language models,
CLIP, surface classification, object detection, GPX.
\end{IEEEkeywords}

% =====================================================================
% PAGE 1 -- Introduction, research approach/methodology, objectives,
% constraints (to be written later, headings only here)
% =====================================================================

\section{Introduction}
% --- NOTE: This section will be written later. ---

\section{Research Approach and Methodology}
% --- NOTE: This section will be written later. ---

\section{Objectives and Constraints}
% --- NOTE: This section will be written later. ---

% =====================================================================
% PAGE 2 -- GROUND DETECTION (Bünyamin Aydemir)
% =====================================================================

\clearpage
\section{Ground Detection: Zero-Shot Classification of the Road Surface}

\emph{Ground detection} determines the ridden surface class for every
video frame and thereby provides the first of the three sub-components
of the bikeability score. Since each frame corresponds to exactly
\emph{one} surface class, the task is formulated as
\emph{whole-image classification} (following the principle
\enquote{whole image $=$ one class}) rather than semantic segmentation.
The input images originate from a camera mounted at the front of the
bicycle handlebar and are sampled at one frame per second.

\subsection{Task Type and Model Class}
Because neither annotated target classes nor a project-specific training
set are available in field operation, the surface is determined
\emph{training-free} (zero-shot). This leads to the choice of
\emph{vision-language models} (VLM) of the CLIP family \cite{clip},
which select the matching class from an image and a set of text prompts
without any fine-tuning. Classical CNNs trained on ImageNet do not know
application-specific classes such as \enquote{gravel path} and are
therefore ruled out.

A distinction that is frequently confused in the current literature is
central here: \emph{linear probing} of self-supervised features
(e.g., DINOv2/DINOv3 \cite{dinov2}) is \emph{not} a zero-shot method.
A linear classifier on frozen features requires trained labels of the
target classes and is therefore \emph{supervised}. Such a method is
deliberately not counted as zero-shot here; at most it could serve as an
explicitly declared \emph{non-zero-shot reference point}.

\subsection{Surface Classes and Prompt Ensembling}
Four cycling-relevant classes are distinguished:
\textbf{Cycleway} (paved bike path), \textbf{Road}
(paved/asphalt road), \textbf{Gravel} (gravel/crushed stone), and
\textbf{Unpaved} (dirt/field track). Since \emph{Cycleway} and
\emph{Road} both consist of smooth asphalt on the surface, the text
prompts deliberately emphasize the \emph{distinguishing context}
(narrow, structurally separated, bicycle pictogram vs. wide, cars,
center line, multiple lanes) rather than only the texture. For each
class, several English descriptions are used, and their text embeddings
are averaged and normalized (\emph{prompt ensembling} as in the CLIP
paper \cite{clip}), which makes the class embeddings more discriminative.
An image is classified via the cosine similarity to these class
embeddings; a softmax yields class probabilities.

\subsection{Compared Models}
To ensure a fair comparison, all models are accessed through the unified
\texttt{open\_clip} API \cite{openclip}. A range from \emph{very light}
to \emph{medium-sized} is deliberately considered -- all local/edge
capable, no closed \enquote{killer model}. Table~\ref{tab:models}
summarizes the considered architectures.

\begin{table}[!t]
\caption{Models compared in the zero-shot benchmark.}
\label{tab:models}
\centering
\small
\begin{tabular}{@{}p{0.55\columnwidth}p{0.35\columnwidth}@{}}
\toprule
\textbf{Model} & \textbf{Characteristic}\\
\midrule
CLIP ViT-B/32, B/16, L/14 \cite{clip} & Classical reference\\
OpenCLIP (LAION-2B, DataComp-XL) \cite{openclip,datacomp} & Dataset influence\\
SigLIP ViT-B/16 \cite{siglip} & Efficient zero-shot standard\\
MobileCLIP S0, S2 \cite{mobileclip} & Latency-optimized\\
EVA02-B/16 \cite{evaclip} & Open zero-shot default\\
\bottomrule
\end{tabular}
\end{table}

\subsection{Metrics}
The benchmark evaluates the models along two axes:
\emph{accuracy} and \emph{latency}.

\paragraph{Accuracy} Since the classes are imbalanced, in addition to
accuracy we mainly report the \emph{balanced accuracy} (mean of the
per-class recalls) and the \emph{macro-F1}. In addition, a
\emph{confusion matrix}, a \emph{bootstrap confidence interval}
(95\,\%, 2000 resamples) on the accuracy, and the \emph{McNemar test}
for pairwise model comparison are computed per model. The McNemar test
checks on the discordant predictions whether an accuracy difference is
\emph{statistically significant} ($p < 0.05$) and not merely noise on a
small test set.

\paragraph{Latency} On fixed hardware, the pure inference latency is
measured at \emph{batch size 1} (real use case: video at the handlebar):
several warm-up runs, followed by $N$ measured single inferences
including GPU synchronization. The mean, standard deviation, the 95th
percentile, and the throughput in frames per second (FPS) are reported.

\subsection{Result Presentation}
The central overall representation is a scatter plot of
\emph{accuracy vs. latency} (Fig.~\ref{fig:pareto}). Each point is a
model (point size $\propto$ number of parameters); the drawn
\emph{Pareto front} connects the models that achieve the best accuracy
for a given latency -- i.e., the best trade-off for use at the bicycle
handlebar. In practice, \textbf{MobileCLIP} variants often lie on the
Pareto front, as they offer a very favorable latency-accuracy ratio
without requiring an oversized model.

\begin{figure}[!t]
\centering
% \includegraphics[width=\columnwidth]{images/pareto_accuracy_latency}
\caption{Accuracy (macro-F1) vs. latency (ms/frame, batch size 1).
Point size $\propto$ number of parameters; the dashed line marks the
Pareto front. \emph{Figure to follow.}}
\label{fig:pareto}
\end{figure}

\begin{figure}[!t]
\centering
% \includegraphics[width=\columnwidth]{images/confusion_matrices}
\caption{Row-normalized confusion matrices per model. They show which
surfaces are confused (typically Gravel~$\leftrightarrow$~Unpaved as
well as Road~$\leftrightarrow$~Cycleway). \emph{Figure to follow.}}
\label{fig:confusion}
\end{figure}

For reproducibility: latency values are strongly hardware-dependent,
therefore the concrete computing device (CPU/GPU model, RAM) as well as
the used \texttt{DEVICE} are documented. All models use the same full
image, the same prompt ensembling, and the same \texttt{open\_clip}
pipeline; observed differences therefore stem from architecture and
training data, not from the evaluation.

% =====================================================================
% PAGE 3 -- OBJECT DETECTION (not Bünyamin -- heading only)
% =====================================================================

\clearpage
\section{Object Detection}
% --- NOTE: This section will be written later. ---

% =====================================================================
% PAGE 4 -- ENVIRONMENT DETECTION (not Bünyamin -- heading only)
% =====================================================================

\clearpage
\section{Environment Detection}
% --- NOTE: This section will be written later. ---

% =====================================================================
% PAGE 5 -- BIKEABILITY SCORE (Bünyamin Aydemir)
% =====================================================================

\clearpage
\section{Computation of the Bikeability Score}

The bikeability score combines the three sub-components -- surface
(ground detection), environment (environment detection), and objects
(object detection) -- into a single value in the interval $[0, 100]$.
Each model output is first transformed into a \emph{sub-score} in
$[0,1]$; the three sub-scores are then combined in a weighted manner.

\subsection{Overall Formula}
For each evaluated frame, the score $A$ results from the weighted mean
of the three sub-scores:
\begin{equation}
\label{eq:score}
A = 100 \cdot \Big( w_g \, S_{\text{ground}}
                  + w_e \, S_{\text{env}}
                  + w_o \, S_{\text{object}} \Big),
\end{equation}
with the weights
\begin{equation}
\label{eq:weights}
w_g = 0.20, \quad w_e = 0.45, \quad w_o = 0.35,
\quad w_g + w_e + w_o = 1.
\end{equation}
Since all sub-scores lie in $[0,1]$ and the weights sum to $1$, the
score is \emph{guaranteed} to be bounded to $[0, 100]$.

\subsection{Ground and Environment Sub-Score}
The surface and environment sub-scores result from a direct lookup of
the respective class quality. If, instead of a single class, a share
distribution $\{(k, p_k)\}$ is present, the weighted mean of the class
qualities is formed:
\begin{equation}
\label{eq:sub_lookup}
S_{\text{ground}} = \sum_{k} q^{g}_{k}\, p_{k},
\qquad
S_{\text{env}} = \sum_{k} q^{e}_{k}\, p_{k},
\end{equation}
where $p_k$ is the share of class $k$ ($\sum_k p_k = 1$) and
$q^{g}_{k}, q^{e}_{k} \in [0,1]$ are the stored quality values.
Table~\ref{tab:quality} lists the (heuristic) quality values used.

\begin{table}[!b]
\caption{Class quality values for the ground and environment sub-score
as well as object penalty factors.}
\label{tab:quality}
\centering
\begin{tabular}{@{}llc@{}}
\toprule
\textbf{Component} & \textbf{Class} & \textbf{Value}\\
\midrule
\multirow{4}{*}{Ground $q^{g}_{k}$}
  & Cycleway & $1.0$\\
  & Road     & $0.7$\\
  & Gravel   & $0.3$\\
  & Unpaved  & $0.1$\\
\midrule
\multirow{3}{*}{Environment $q^{e}_{k}$}
  & Vegetation & $1.0$\\
  & Water      & $1.0$\\
  & City       & $0.4$\\
\midrule
\multirow{3}{*}{Object penalty $\lambda_k$}
  & Car     & $+1.0$\\
  & Bus     & $+2.0$\\
  & Bicycle & $-0.2$\\
\bottomrule
\end{tabular}
\end{table}

\subsection{Object Sub-Score}
The object sub-score evaluates the detected object counts $n_k$
(e.g., cars, buses, cyclists). Via a \emph{saturation function}, a
weighted object density is mapped onto $(0, 1]$:
\begin{equation}
\label{eq:sub_object}
S_{\text{object}} = \exp\!\Big( -\max\big(\rho, 0\big) \Big),
\qquad
\rho = \frac{1}{N}\sum_{k} \lambda_k\, n_k,
\end{equation}
with the number of evaluated frames $N$ and the penalty factors
$\lambda_k$ from Table~\ref{tab:quality}. Positive penalty factors
(e.g., cars, buses) lower the sub-score, while bicycle-friendly objects
such as other cyclists receive a \emph{negative} factor and raise the
score slightly. The exponential form is robust against outliers: a
single overcrowded frame does not dominate the overall score.

\subsection{Route Evaluation and Aggregation}
To evaluate a complete ride, the video is sampled at fixed intervals
(every $5$\,s, i.e., every $\lfloor \text{FPS}\cdot 5\rfloor$-th frame)
and a score is computed for each frame according to
Eq.~\eqref{eq:score}. Each row is stored in a CSV with a timestamp
(seconds from video start, \texttt{HH:MM:SS}, and optionally an absolute
ISO timestamp). Via the absolute timestamp, the frame scores can then be
\emph{merged} with the GPX track points of the ride and located
geographically. The \emph{average} bikeability score of the route
results as the mean over all $M$ evaluated frames:
\begin{equation}
\label{eq:avg}
\bar{A} = \frac{1}{M} \sum_{i=1}^{M} A_i .
\end{equation}

Figure~\ref{fig:route} shows the final result: the bikeability score
colored along the GPX route, which visualizes the section-wise
cycling-friendliness of the ridden route.

\begin{figure}[!t]
\centering
% \includegraphics[width=\columnwidth]{images/bikeability_route_map}
\caption{Bikeability score colored along the GPX route. The frame scores
are merged with the GPX track points via the absolute timestamp.
\emph{Figure to follow.}}
\label{fig:route}
\end{figure}

% =====================================================================
% CONCLUSION (to be written later, heading only)
% =====================================================================

\section{Conclusion}
% --- NOTE: This section will be written later. ---

% =====================================================================
% Bibliography
% =====================================================================

\bibliographystyle{IEEEtran}
\bibliography{references}

\end{document}
